% !TEX root = ../main.tex
% =====================================================================
% Experiments section --- CONFERENCE (CVPR) version, written per
% docs/reviews/paper_outline_cvpr.md §3 (compact, contribution-driven).
% Same canonical numbers and honesty rules as the dissertation Ch.5
% (dissertation/sec/evaluation.tex); prose is genre-adapted, NOT identical.
% Sources: project_review_202607.md §1.4, benchmark_mmbv_final_analysis.md,
% benchmark_mmbv_final_gpt4judge.md (+_official_agg), dense 16f/32f reports,
% benchmark_v2_agqa.md, progress.md §14.1/14.2/16.1/16.2.
% Judge discipline: paper-level & dimension-level = gpt-4-turbo; runs=3 ±std
% noise argument = qwen-max; duration/attribution derived from qwen-max raw.
% AGQA + annotation audit: short paragraphs here, tables in supplementary.
% =====================================================================

\section{Experiments}
\label{sec:eval}

\subsection{Setup}
\label{subsec:exp-setup}

\textbf{Benchmark and protocol.} We evaluate on a 150-question stratified subset
(seed 42; 136 videos) of MMBench-Video~\cite{fang2024mmbenchvideo}, oversampling
the Temporal Reasoning (TR) and Hallucination (HL) dimensions; scores are therefore
\emph{not comparable to the public leaderboard} (re-weighting by the official
distribution moves totals by $\leq 0.04$). We replicate the official
VLMEvalKit~\cite{duan2024vlmevalkit} judging protocol (0--3 LLM judge, official
multi-label dimension aggregation) under two judges: \texttt{qwen-max} as-run and
a full \texttt{gpt-4-turbo}~\cite{openai2023gpt4} re-judge (1{,}350 judgements).
Unless noted, dimension-level numbers use gpt-4-turbo; the two settings agree
within 0.015 on all totals (per-question agreement 0.76--0.81). Duration and
attribution analyses derive from qwen-max raw scores.

\textbf{Systems.} All systems share the same LLM
(\texttt{qwen-plus-latest}~\cite{bai2023qwen}) and VLM
(\texttt{qwen-vl-plus}~\cite{bai2023qwenvl}): our agent;
\textbf{vlm\_direct@8} (8 uniform frames); \textbf{vlm\_transcript@8}, the
\emph{fair baseline} that additionally receives the full ASR transcript, holding
modalities equal to the agent's; and the v1 scene-graph pipeline
(\cref{sec:diagnosis}). Main systems are run 3 times (mean$\pm$std); all token,
call, and frame counts are API-reported, never estimated.

\subsection{Main Results and Attribution}
\label{subsec:exp-main}

\begin{table}[t]
  \centering
  \footnotesize
  \caption{Overall results (0--3, mean$\pm$std over 3 runs) and per-question cost.
  v1: legacy scene-graph pipeline, single run.}
  \label{tab:main}
  \setlength{\tabcolsep}{2pt}
  \begin{tabular}{@{}lcccc@{}}
    \toprule
    Method & gpt-4-turbo & qwen-max & Frames/Q & Tokens/Q \\
    \midrule
    \textbf{agent (ours)} & $\mathbf{1.978{\pm}0.097}$ & $\mathbf{1.984{\pm}0.101}$ & \textbf{3.6} & 16{,}655 \\
    vlm\_transcript@8 & $1.713{\pm}0.019$ & $1.727{\pm}0.020$ & 8.0 & 7{,}299 \\
    vlm\_direct@8 & $1.491{\pm}0.019$ & $1.478{\pm}0.025$ & 8.0 & 6{,}748 \\
    v1 pipeline & --- & 1.193 & --- & --- \\
    \bottomrule
  \end{tabular}
\end{table}

\Cref{tab:main}: the agent leads the fair baseline by 0.265 (gpt-4-turbo; 0.257
under qwen-max), and the reversal survives run noise --- the qwen-max gap exceeds
the sum of both standard deviations ($0.257 > 0.101 + 0.020$). The agent's larger
variance reflects exploration-path randomness across runs. v1's 1.193 --- below
even the direct baseline --- shows that a scene graph used as \emph{answer source}
is a liability; the gain comes from using it as an index
(\cref{subsubsec:provenance}).

\textbf{Attribution.} The fair baseline splits the total gain cleanly in two
(qwen-max): adding the ASR modality to the direct baseline is worth $+0.249$
($1.478 \to 1.727$); the architecture --- lazy evidence memory, provenance
pull-back, confidence-driven exploration --- adds $+0.257$ on top
($1.727 \to 1.984$; $+0.265$ under gpt-4-turbo) with modalities held fixed. The
agent does not win merely by hearing the narration.

\subsection{Analysis}
\label{subsec:exp-analysis}

\begin{table*}[t]
  \centering
  \small
  \caption{Per-dimension results, MMBench-Video Table-3 layout (gpt-4-turbo judge,
  official multi-label aggregation, 0--3). Full tables with std for both judges in
  the supplementary material.}
  \label{tab:dimensions}
  \begin{tabular}{@{}lcccccccccccc@{}}
    \toprule
    Model & Overall & CP & FP-S & FP-C & HL & \emph{P. Mean} & LR & AR & RR & CSR & TR & \emph{R. Mean} \\
    \midrule
    agent (ours) & \textbf{1.98} & \textbf{2.18} & \textbf{2.08} & \textbf{1.47} & \textbf{2.29} & \textbf{2.04} & \textbf{2.18} & 2.07 & 2.07 & 2.21 & 1.74 & 1.98 \\
    vlm\_transcript & 1.71 & 1.78 & 1.81 & 0.90 & 0.78 & 1.42 & 1.73 & \textbf{2.50} & \textbf{2.12} & \textbf{2.28} & \textbf{1.87} & \textbf{2.07} \\
    vlm\_direct & 1.49 & 1.53 & 1.64 & 0.77 & 0.64 & 1.28 & 1.39 & 2.21 & 1.81 & 2.26 & 1.52 & 1.79 \\
    \bottomrule
  \end{tabular}
\end{table*}

\begin{table}[t]
  \centering
  \small
  \caption{Overall by video duration (qwen-max raw scores, runs=3 means).}
  \label{tab:duration}
  \begin{tabular}{@{}lccc@{}}
    \toprule
    Duration & agent & vlm\_transcript & vlm\_direct \\
    \midrule
    $<$90\,s \hfill ($n{=}52$) & 1.81 & 1.76 & 1.57 \\
    90--180\,s \hfill ($n{=}50$) & \textbf{2.10} & 1.55 & 1.34 \\
    $>$180\,s \hfill ($n{=}48$) & \textbf{2.05} & 1.87 & 1.52 \\
    \bottomrule
  \end{tabular}
\end{table}

\textbf{Where the gain lives.} By duration (\cref{tab:duration}), the agent and
the fair baseline are on par below 90\,s (1.81 vs.\ 1.76) and separate sharply
beyond it (2.10 vs.\ 1.55 at 90--180\,s; 2.05 vs.\ 1.87 above 180\,s): uniform
sampling dilutes with length, guided sampling does not. By dimension
(\cref{tab:dimensions}), the win is a \emph{perception} win (P.\ Mean 2.04 vs.\
1.42) while reasoning stays near parity (1.98 vs.\ 2.07) --- all systems share
one LLM, so the architecture buys seeing accurately, not thinking harder.
Hallucination resistance is the standout: HL 2.29 vs.\ 0.78
($\approx$2.9$\times$; 2.42 vs.\ 1.04 under qwen-max), which we attribute to the
grounding rule that forbids answering ``no'' about unexplored content
(\cref{subsubsec:grounding}). We disclose the one losing dimension: TR (1.74 vs.\
1.87). The baseline's own TR jump with ASR (1.52 $\to$ 1.87) shows the temporal
signal is largely \emph{modality}-carried, which is why we claim evidence
localization and integration, not improved temporal reasoning.

\textbf{Qualitative example.} Asked \emph{``Did NIO rank the 20th place in the
chart?''} (HL item mmbv-1185), free retrieval finds a company list without NIO; instead
of answering ``no'' directly, the agent explores 120--150\,s and then 60--120\,s
(12 frames total), verifies that chart blocks carry explicit rank numbers, reasons
that the explored range must contain rank \#20, and only then answers ``no'' with
the evidence trail --- a negative claim earned by targeted perception (full trace
in the supplementary material).

\subsection{Frame Scaling}
\label{subsec:exp-scaling}

Can the baselines catch up by simply looking at more frames? We densify both
baselines to 16 and 32 uniform frames (single run, both judges).
Three questions are excluded from \emph{all} conditions (common-147): dense
sampling deterministically triggers the hosting provider's content moderation
(errors grow 0/1/3 at 8/16/32 frames, 100\% reproducible; the agent and 8f
baselines are unaffected), so keeping them would credit the agent for
infrastructure failures.\footnote{The full-150 accounting, which scores the blocked
items as 0 and thus \emph{underestimates} the dense baselines, is in the
supplementary material; conclusions are unchanged.}

\begin{figure}[t]
  \centering
  % Figure is a PDF pre-converted from fig/frame_scaling.svg (rsvg-convert -f pdf).
  % Using the PDF avoids invoking Inkscape on every compile: much faster on
  % Overleaf, and it works in any TeX install without shell-escape.
  % SVG fallback: keep \usepackage{svg} and swap this line back to
  %   \includesvg[inkscapelatex=false,width=\linewidth]{fig/frame_scaling}
  \includegraphics[width=\linewidth]{fig/frame_scaling}
  \caption{Frame-accuracy scaling (common-147, trial-1): dense uniform sampling
  never overtakes guided sampling under either judge.}
  \label{fig:framescaling}
\end{figure}

\begin{table}[t]
  \centering
  \small
  \caption{Frame-scaling (0--3, common-147, trial-1). agent = its
  $\sim$3.6-frame operating point; $\Delta$32f = gap remaining at 32 frames.}
  \label{tab:framescaling}
  \setlength{\tabcolsep}{4pt}
  \begin{tabular}{@{}lccccc@{}}
    \toprule
    Method (judge) & 8f & 16f & 32f & agent & $\Delta$32f \\
    \midrule
    direct (qwen) & 1.442 & 1.497 & 1.680 & \textbf{2.109} & $-0.429$ \\
    transcript (qwen) & 1.735 & 1.660 & 1.823 & \textbf{2.109} & $-0.286$ \\
    direct (gpt-4-t) & 1.510 & 1.707 & 1.857 & \textbf{2.088} & $-0.231$ \\
    transcript (gpt-4-t) & 1.694 & 1.871 & 1.932 & \textbf{2.088} & $-0.156$ \\
    \bottomrule
  \end{tabular}
\end{table}

\Cref{fig:framescaling,tab:framescaling}: vlm\_direct gains only $+0.12$--$0.17$
per frame doubling and at 32f still trails by 0.23--0.43; log-linear extrapolation
puts break-even at $\sim$81 (gpt-4-turbo) to $\sim$389 (qwen-max) frames ---
\textbf{22--108$\times$ the agent's frame budget}. The transcript baseline at 32f
reaches the agent's noise band only under the more lenient gpt-4-turbo judge
(1.932 vs.\ 1.957, gap 0.025) --- at the price of the full narration modality plus
$\sim$9$\times$ the frames, and still 0.29 behind under the stricter judge. Dense
sampling thus sits on a strictly worse cost--accuracy frontier, and grows more
fragile with density (the moderation failure above).

\subsection{Negative Results and Limitations}
\label{subsec:exp-limits}

\textbf{Falsified alternatives.} (i)~\emph{Oracle routing}: zero-exploration
questions score higher than escalated ones (2.259 vs.\ 1.551, single-run
diagnostic), suggesting a router that skips exploration per dimension. Routing on
\emph{gold} labels --- an upper bound for any learned router --- drops the total
from 1.933 to 1.733 and hurts every routed dimension (TR $-0.37$, FP-C $-0.27$,
LR $-0.40$; matched questions: 2 improved, 13 worsened): the correlation is
selection bias (easy questions do not trigger exploration), and the direction is
closed. (ii)~\emph{Deeper evidence}: raising the per-window frame cap ($6 \to 9$)
with a restructured caption prompt recovered 0 of 4 targeted failures --- the
residual gap is model capability, not evidence depth. The dominant remaining
failure class is localization on videos over 200\,s.

\textbf{Transfer and annotation quality.} On a 70-question English AGQA
subset~\cite{grundemclaughlin2021agqa} (single run), the agent's profile
transfers: overall 0.436 (LLM-judge), strongest on duration questions (0.682),
weak on open-vocabulary naming (0.190) and sequencing (0.286); per-category tables and an
analysis of exact-match scoring artefacts are in the supplementary material. An
audit of 30 sampled gold answers found 97\% fully evidence-supported (score
ceiling $\approx$0.98), so the remaining headroom reflects capability, not label
noise (audit protocol and its LLM-proxy caveat in the supplementary material).

\textbf{Limitations.} Results rest on one benchmark subset (not
leaderboard-comparable; composition effect $\leq 0.04$); judge self-preference is
refuted by the cross-judge agreement above; the agent costs $\approx$2.1$\times$
the baseline's tokens and $\approx$6$\times$ its latency --- the contribution is
frame efficiency and long-video accuracy, not cost; backbone models are hosted
behind rolling aliases (all runs within June--July 2026).
